\chapter{Safety}
\label{ch:safety}

\section*{What this chapter is for}

\reported{Safety questions are moderately common and role-dependent, and appear
more often as a named failure mode than as a dedicated
round.}{field-guide} Where they do appear, they are usually the same question
in different clothes: what can this system be made to do that you did not
intend?

\section*{The questions}

\begin{bankq}{What is prompt injection, and why is it worse in an agent?}
\qasked
Sometimes named directly, more often reached through a scenario about a
system that did something it should not have.

\qtested
Whether you locate the problem at the data boundary or treat it as a
prompt-writing problem.

\qstaff
The system cannot distinguish instructions from data, because both are text in
the same channel. Anything that reaches the context can be read as an
instruction, and that includes content the system fetched itself \dash{} a
retrieved document, a tool result, a web page, a file.

That last part is what makes it different from a user typing something
adversarial, and it is the part people miss: the attack does not have to come
from the person talking to the system. \reported{The amplification in agentic
systems is that the model can act, so a successful injection is no longer a
matter of what it says but of what it does with the tools it
holds.}{agentic-injection}

The answer is the same one as Chapter~\ref{ch:agents}: controls live at the
tool boundary. An injected instruction that persuades the model to call a write
tool still meets a tool that refuses without an approval bound to what a person
saw. \judgement{Filtering inputs for suspicious phrasing is worth doing and is
not a control \dash{} it is a filter with an unknown false-negative rate
protecting something that should not have needed protecting.}

There is a second-order point worth making if there is room: a system that
refuses to function whenever it sees something suspicious has been
denial-of-serviced by its own defences. The requirement is to keep working
\emph{and} not act.

\qsenior
A correct description of a user injecting instructions, answered with input
filtering and a firmer system prompt. It defends the channel that is easiest to
defend and leaves the retrieved-content channel open.

\qcontested
Whether the problem is solvable at the model layer. \judgement{The defensible
position is that it is not, currently, and that architectures should assume
injection succeeds and be designed so that succeeding is not
useful.}\source{mcp-prompt-injection}

\qdrill
For a system you know, list every path by which text you do not control reaches
the context. Most people find one more than they expected.
\end{bankq}

\begin{bankq}{What is the worst thing this system can do without a person?}
\qasked
A design question, and one of the few in this book that is better asked of
yourself than waited for.

\qtested
Whether you can state a blast radius. It is a Staff question in a safety
costume: it is really about whether you know what you built.

\qstaff
Answer it literally, as one sentence, then say what bounds it. The bound is
the union of what every tool can do with the credentials it holds \dash{} not
what the model is likely to do, which is a different and unanswerable question.

Then the reductions, in order of effectiveness: fewer tools that write; write
tools that require an approval bound to a specific action; credentials scoped
per tool rather than one capable credential shared; and anything irreversible
either behind a person or made reversible.

\judgement{The question is a good one to ask yourself about anything you have
built, and being unable to answer it in a sentence is itself the answer.}

\qsenior
Discussion of guardrails, content filtering and monitoring. Those are about
what the model produces; the question was about what the system can do.

\qcontested
How much autonomy is acceptable for reversible actions varies entirely by
domain, and no general answer exists.

\qdrill
Write the sentence for a system you have built. If it takes more than one
sentence, the tools are too capable.
\end{bankq}

\begin{bankq}{How do you stop it leaking things it should not?}
\qasked
Usually framed around a specific worry the team already has, such as one
user seeing another user's documents.

\qtested
Whether you think of this as an output problem or an input problem.

\qstaff
Mostly it is an input problem. Data that never entered the context cannot be
returned from it, so the first control is retrieval that respects the asking
user's permissions \dash{} filtering at query time rather than filtering the
answer afterwards.

Then the output side, which catches what the first misses: internal
identifiers, system details and permission names should not appear in
user-facing text. That last one is easy to miss \dash{} a refusal saying
\emph{you lack the inventory-write permission} satisfies a naive check for
refusing and leaks the shape of your authorisation model in the same sentence.

And separately: keep credentials and remote error text out of traces. A trace
goes to a collector and is read by people who did not need to see either.

\qsenior
Output filtering and redaction. Necessary, and it is the last line rather than
the first.

\qcontested
Little in principle. In practice, whether permission-aware retrieval is
feasible depends on whether the corpus carries usable permissions, and
frequently it does not.

\qdrill
Take a refusal message from a system you know and check whether it names
anything internal. It very often does.
\end{bankq}

\section*{What this chapter does not claim}

The frequency claim comes from one corpus. The injection material is consistent
with published practitioner writing on agentic amplification and tool blast
radius, cited where used.

The position that injection is not currently solvable at the model layer is a
judgement, and it is one on which reasonable people differ.
